\subsubsection{BackEnd}
Mit Beginn der Backend-Entwicklung wurden in erster Linie die Grundstrukturen implementiert, um der \textit{MVVM}-Architektur gerecht zu werden. Dabei lag 
der Fokus bei der Instanziierung der notwendigen Komponenten der \textit{Android Architecture Components}. Nach chronologischer Reihenfolge wurden die Klassen 
erstellt und mit den dazugehörigen Funktionen und Methoden versehen. Angefangen mit der Entity-Klasse zur Beschreibung des Objekts und deren allgemeiner Aufbau, 
dieser bereits in der Konzeption (\ref{chap:Konzeption}) unter dem Datenmodell (\ref{chap:Datenmodell}) festgehalten wurde. Darauffolgend wurde das „Data Object“ 
erstellt, welches die Zugriffe der Datenbankobjekte managed. Abschließend im Bereich der Datenbank ein Room-Layer, um die eigentliche Datenbank zu 
instanziieren. 
\\ 
Zur Modularisierung und zur Generierung einer Schnittstelle für den Bezug zu mehreren Datenbeziehungspunkten wurde ein Repository erstellt, das eine saubere 
\acs{API} für den Datenzugriff auf den Rest der Anwendung bietet. Um die vorhandenen Klassen zum Datentransfer und zur Persistierung der Daten mit der 
eigentlichen Benutzeroberfläche zu verbinden, wurde ein ViewModel implementiert, welches die Daten zwischen den einzelnen Komponenten teilt und bereitstellt. 
\\ 
\linebreak 
Nach Aufzählung der einzelnen Bestandteile wird auf diese nun genauer eingegangen. 
\\ 
\linebreak
In einer Java-Klasse wird mithilfe der gegebenen Bibliothek „Room“ ein Entity-Objket erstellt. Hierbei gibt es eine eindeutige Annotation, die die Klasse und 
deren beinhalteten Variablen deklariert. So wird, dem folgend aufgeführten Code-Beispiel (\ref{code:entity}) zu entnehmen, eine Klasse in eine Entity und somit in eine 
Datenbank-Tabelle konvertiert. In dieser Tabelle, definiert als \textit{„object\_table“}, gibt es weitere Variablen, die mit Annotationen versehen sind. Damit werden diese 
gemäß der Anforderungen des Konzepts definiert. Diese Variablen repräsentieren die Attribute des Datenbankschematas und stellen die einzelnen Informations- , bzw. 
-Datenbankspalten dar. Zur Veranschaulichung, die Initialisierung der \acs{ID}-Vergabe eines Objekts. Diese wird als \textit{„id“}-Spalte und ebenso als 
Primärschlüssel\footnote{Einmaliger und eindeutiger Wert einer Tabelle, bzw. eines Attributs, um dieses eindeutig zu kennzeichnen.}-Variable deklariert.
\begin{lstlisting}[language=C,
    frame=single,           % Ein Rahmen um den Code
    framexleftmargin=-1pt,  % Rahmen link von den Zahlen
    style=algoBericht,
    label={code:entity},
    captionpos=b,           % Caption unter den Code setzen
caption={Entity Code zur Initialisierung der Objekte}]
@Entity(tableName = "object_table")
public class Object {

    @ColumnInfo(name = "id")
    @NonNull
    @PrimaryKey(autoGenerate = true)
    private int id;

    public int getId() { return this.id; }
    public void setId(int id) { this.id = id; }

    ... 
}
\end{lstlisting}
\pagebreak
Das mit der Entity-klasse kommunizierende Modul ist das \textit{„Data Object“}, welches als Interface angelegt wurde. Das „Object Dao“ verwendet ebenso die Bibliothek 
„Room“ und wandelt die Java-Klasse per Annotation in ein „Dao“ um, dieses beinhaltet hauptsächlich die SQL-Queries zur Datenabfrage. 
\begin{lstlisting}[language=C,
    frame=single,           % Ein Rahmen um den Code
    framexleftmargin=-1pt,  % Rahmen link von den Zahlen
    style=algoBericht,
    label={code:query},
    captionpos=b,           % Caption unter den Code setzen
caption={SQL-Query zur Abfrage der Objekt-Namen}]
@Query("SELECT * FROM object_table ORDER BY name")
LiveData<List<Object>> getObjectName();
\end{lstlisting}
Anschließend nachdem die beiden Klassen erstellt waren, wurde der Datenbank-Layer auf der eigentlichen Datenbank implementiert. Über einen „Builder“ wird die 
Datenbank-Instanz erzeugt und mit einem Namen versehen. Im Fall des Assistenzsystems als \textit{„object\_database“} deklariert. 
\begin{lstlisting}[language=C,
    frame=single,           % Ein Rahmen um den Code
    framexleftmargin=-1pt,  % Rahmen link von den Zahlen
    style=algoBericht,
    label={code:dblayer},
    captionpos=b,           % Caption unter den Code setzen
caption={Erzeugung des Datenbank-Layers „Room“}]
@Database(entities = {Object.class}, version = 1, exportSchema = false)
public abstract class ObjectRoomDatabase extends RoomDatabase {
    ...

    INSTANCE = Room.databaseBuilder(context.getApplicationContext(),
    ObjectRoomDatabase.class, "object_database")
    .addCallback(sRoomDatabaseCallback)
    .build();

    ...
}
\end{lstlisting}
Letzter wichtiger zu implementierender Aspekt war das Kommunikations-Modul zwischen den Datenbank-Transfers und der Benutzeroberfläche, das ViewModel. 
In dieser Klasse wird eine Liste erstellt, welche den Wert des zu speichernden Objekts besitzt. Diese werden über eine „get“-Funktion aus dem zuvor 
instanziierten Repository geladen und in der erzeugten Liste lokal abgelegt. Diese Liste wird dann für die Präsentation der Informationen auf der 
Nutzeroberfläche verwendet. 
\begin{figure}[hbt!]
    \centering
    \includegraphics[width=5cm,height=5cm,keepaspectratio]{4Umsetzung/Bilder/cjt_logo_tracking.png}
    \caption{Marker zur Erkennung der Ausgangsposition}
    \label{pic:initialMarker}
\end{figure}
